\documentclass[12pt]{article}
\usepackage[utf8]{inputenc}
\usepackage[T1]{fontenc}
\usepackage{hyperref}
\title{Bunu, Jiongnai Language Resources for K-12 Curriculum Development}
\author{Compiled by Claude Sonnet 5.5 (claude-sonnet-5-5) for Andr\'as Kornai}
\date{}
\begin{document}
\maketitle

\section{Introduction}
Jiongnai (Kiong Nai; ISO 639-3 pnu) is a Hmongic language of the Hmong-Mien family, spoken in Jinxiu County, Guangxi, China, by about 1,100 people (1999 figure cited by Wikipedia). Speakers are officially classified as Yao. Sources differ on classification: some place it with Bunu, others treat it as an independent branch of Hmongic. Two dialects, Longhua and Liuxiang, are reported. The search found no educational materials; only general reference pages.

\section{Primary Educational Resources}
None were found in the search results.

\section{Linguistic and Reference Materials}
\begin{itemize}
\item Wikipedia, Kiong Nai language (overview, dialects, classification, speaker figures): \url{https://en.wikipedia.org/wiki/Kiong_Nai_language}
\item Wikipedia, Bu--Nao languages (context for the Bunu label): \url{https://en.wikipedia.org/wiki/Bu%E2%80%93Nao_languages}
\item Verbix language page for Bunu, Jiongnai: \url{https://docs.verbix.com/Languages/BunuJiongnai}
\item Austronesian Basic Vocabulary Database page for Jiongnai (Longhua), a comparative word list: \url{https://abvd.eva.mpg.de/austronesian/language.php?id=768}
\item ASJP entry (Jinxiu Changtong Yao), word list data: \url{https://asjp.clld.org/languages/YAO_GUANGXI_JINXIU_CHANGTONG}
\end{itemize}

\section{Religious and Community Resources}
None were found.

\section{Digital and Audio Resources}
No audio, video or corpus resources were found beyond the word lists above.

\section{Curriculum Development Context}
Very little exists for this language. Wikipedia cites Mao and Li (2002) as a source on the dialects, but the search did not locate that work online. A curriculum would require fresh documentation with the community, a practical orthography decision, and attention to Chinese-medium schooling in Jinxiu County. Chinese-language sources should be searched in further work.
\end{document}